\begin{remark}[One simultaneous admissible rate regime]
\label{rem:v23-compatible-rates}
The scaling assumptions used by the endpoint, count--endpoint and stopped
transfer theorems have a common nonempty regime.  For example, let
\[
 \delta_n=n^{-1},\qquad
 k_n=\left\lfloor\frac{n^2}{\log n}\right\rfloor,
 \qquad
 j_n=2\left\lceil C\log n\right\rceil
\]
with \(C>|\log\tau|^{-1}\).  Then
\[
 k_n\delta_n^2\log(1/\delta_n)\to1,\qquad
 j_n\delta_n\to0,\qquad
 k_n\tau^{j_n}\to0,
\]
and also
\[
 j_n=o(\sqrt{k_n}),
 \qquad
 \frac{\log(j_n\sqrt{k_n})}{j_n^2}\to0.
\]
Thus the stated physical and statistical asymptotics do not rely on an empty
intersection of rate conditions.
\end{remark}

\begin{remark}[Exact source of the cap success-ratio estimate]
\label{rem:v23-cap-source}
The success-law input used in Lemma~\ref{lem:v22-reference-cap} is the selected
channel summand in Theorem~\ref{thm:g-stability}, equation
\eqref{eq:g-competition}, evaluated at the fixed physical window.  That
formula gives the factor \(\sinh(j\gamma)^{-1}\), the quadratic offset factor
and a smooth remainder with uniform fixed-order parameter derivatives.  On a
compact local parameter set, Taylor expansion of this exact formula gives the
prefactor ratio \(1+O_K(\delta_n)\) and the exponent variation
\(j_nO_K(\delta_n)\) used in
\eqref{eq:v22-success-ratio}.  Hence \(j_n\delta_n\to0\) yields the required
uniform ratio \(1+o(1)\) for the single reference-based cap sequence.
\end{remark}
